\DIFbluesubsection{LLM Fallback}
\label{sec:llm-fallback}

\DIFaddbegin
Algorithm~\ref{alg:funspec-gen} presents the idealized symbolic-execution-driven path and intentionally omits the fallback layer; the LLM-fallback layer wraps that algorithm and is triggered in exactly two cases. (i) When the target function is self-recursive or manipulates recursive data structures: the symbolic executor neither computes recursive fixpoints nor enumerates the unbounded cell set of a recursive shape, so the fallback is invoked directly at the function entry instead of running the algorithm. (ii) When $\textit{SP}(\cdot)$ at the function exit returns $\textit{SymExecFailed}$: this covers the residual cases in which symbolic execution cannot produce a usable summary, including external calls, path explosion, and runtime errors. Loops and ordinary callees do not enter fallback; they are handled by $\mathrm{LoopInvGen}$ and recursive $\mathrm{FuncSpec}$ respectively. Their residual verifier failures are sent to the verifier-guided refinement layer, which also wraps Algorithm~\ref{alg:funspec-gen} rather than being inlined.
\DIFaddend

\DIFaddbegin
In the fallback, the LLM acts as a semantic summarization component for the missing symbolic‑execution step. It receives the current code fragment, available symbolic states and path conditions, previously generated callee specifications or loop summaries, and an ACSL skeleton. The LLM is tasked with following the relevant paths, tracking changes to variables and memory locations, and proposing ACSL clauses that summarize the observed behavior. The resulting candidates must still be proved by Frama‑C/WP. 
\DIFaddend

%The fallback uses the LLM as a semantic summarization component for the missing symbolic-execution step, proposing ACSL candidates that must still be proved by Frama-C/WP. It receives the current code fragment, available symbolic states and path conditions, previously generated callee specifications or loop summaries, and an ACSL skeleton, and is instructed to follow the relevant paths, to track changes to variables and memory locations, and to propose ACSL clauses that summarize the observed behavior.

\DIFaddbegin
We use this fallback for four function-level tasks. \emph{Precondition generation} asks the LLM to infer candidate \texttt{requires} clauses describing the validity, separation and bound constraints. \emph{Assigns-clause generation} asks the LLM to infer the function's footprint, i.e., \texttt{assigns} clauses describing which visible memory locations may be modified. \emph{Postcondition generation} asks the LLM to summarize the relation between the pre-state and the post-state, filling candidate \texttt{ensures} clauses. \emph{Inductive-predicate generation} asks the LLM to decide whether the program requires an inductive predicate and, when needed, to propose one for capturing inductive relations in recursive data structures or recursive functions.
\DIFaddend

\DIFaddbegin
Fallback candidates are accepted only when Frama-C/WP proves them under the configured memory model and solver settings; verifier failures are handled by the refinement loop in Section~\ref{sec:refinement} (Table~\ref{tab:refine-strategy}). Fallback can therefore cover cases beyond the exact symbolic-execution fragment but 
unlike the postconditions computed by \(\textit{SP}\) (Proposition~\ref{thm:sp-sound-complete}), the results produced by the fallback are sound over‑approximations verified by the prover.
\DIFaddend

%does not extend the strongest-postcondition guarantee of Proposition~\ref{thm:sp-sound-complete}.
